We know that for every regular scheme having an ample sheaf, the Grothendieck groups formed with locally free sheaves of finite type and with coherent sheaves are identical. In particular, this is the case for @@M0@@.

Let @@M1@@ be a coherent @@M2@@-Module supported on @@M3@@; then @@M4@@ is a perfect @@M5@@-Module. Indeed, it is perfect at the points of @@M6@@, since @@M7@@ is regular at these points, and at the points of @@M8@@ since it vanishes there.

If @@M9@@ is the class of @@M10@@ in @@M11@@, then @@M12@@ belongs to the image of @@M13@@. Indeed, we can filter @@M14@@ by the @@M15@@. This filtration is finite, since, @@M16@@ being noetherian, there is an integer @@M17@@ such that @@M18@@; the quotients of this filtration are then @@M19@@-Modules, so their classes in @@M20@@ belong to the image of @@M21@@. Since @@M22@@ is the sum of the classes of these quotients, @@M23@@ also belongs to this image.

Let now @@M24@@. We can find a perfect complex @@M25@@ whose class is @@M26@@. Consider the distinguished triangle:
@@M27@@
where the horizontal arrow is the canonical homomorphism. This arrow is an isomorphism on the complement of @@M28@@, and consequently the cohomology sheaves of @@M29@@ are coherent sheaves supported on @@M30@@. By the preceding argument, these are perfect @@M31@@-Modules, and their classes belong to the image of @@M32@@. Thus we have:
@@M33@@
and consequently the homomorphism @@M34@@ is surjective.



% Macros for Exposé VII continuation
\providecommand{\Q}{\mathbb Q}
\providecommand{\bfP}{\mathbf P}
\providecommand{\gr}{\operatorname{gr}}
\providecommand{\Id}{\operatorname{Id}}
\providecommand{\Fil}{\operatorname{Fil}}
\providecommand{\Imf}{\operatorname{Im}}
\providecommand{\rank}{\operatorname{rank}}
\providecommand{\codim}{\operatorname{codim}}
\providecommand{\cJ}{\mathcal J}
\providecommand{\cI}{\mathcal I}
\providecommand{\cN}{\mathcal N}
\providecommand{\cE}{\mathcal E}
\providecommand{\cO}{\mathcal O}


